\FigureLayoutDeclare{img/2013/shanghai_spring/q25_fig1.png}{4.0cm}{0bp 0bp 0bp 0bp}

\chapter{2013年上海卷（春）}

\section{填空题}

\begin{problem}
函数 \( y = \log_{2}(x + 2) \) 的定义域是\fillinblank{}.
\end{problem}

\begin{answer}
\( (-2, +\infty) \)
\end{answer}

\begin{solution}
对数的真数必须为正，因此必须满足 \(x+2>0\)，解得 \(x>-2\). 所以函数的定义域为 \((-2,+\infty)\).
\end{solution}

\begin{problem}
方程 \( 2^{x}=8 \) 的解是\fillinblank{}.
\end{problem}

\begin{answer}
\(3\)
\end{answer}

\begin{solution}
因为 \(8=2^3\)，所以原方程化为 \(2^x=2^3\). 指数函数 \(2^x\) 在 \(\R\) 上单调递增，故 \(x=3\).
\end{solution}

\begin{problem}
抛物线 \( y^{2}=8x \) 的准线方程是\fillinblank{}.
\end{problem}

\begin{answer}
\(x=-2\)
\end{answer}

\begin{solution}
抛物线的标准方程为 \(y^2=4px\)，其准线方程为 \(x=-p\). 由 \(4p=8\) 得 \(p=2\)，所以准线方程为 \(x=-2\).
\end{solution}

\begin{problem}
函数 \( y = 2 \sin x \) 的最小正周期是\fillinblank{}.
\end{problem}

\begin{answer}
\(2\pi\)
\end{answer}

\begin{solution}
正弦函数 \(\sin x\) 的最小正周期为 \(2\pi\)，前面的系数 \(2\) 只改变函数值的振幅，不改变周期. 因此 \(y=2\sin x\) 的最小正周期为 \(2\pi\).
\end{solution}

\begin{problem}
已知向量 \( \bs{a} = (1, k) \)，\( \bs{b} = (9, k - 6) \). 若 \( \bs{a} \parallel \bs{b} \)，则实数 \( k =\fillinblank{}\).
\end{problem}

\begin{answer}
\(-\frac{3}{4}\)
\end{answer}

\begin{solution}
两向量平行的充要条件是其坐标行列式为零. 因而
\[
\begin{vmatrix}1&k\\9&k-6\end{vmatrix}=1\cdot(k-6)-9k=0
\]
解得 \(-8k-6=0\)，即 \(k=-\frac34\). 两向量的第一坐标均非零，不存在零向量的额外情形.
\end{solution}

\begin{problem}
函数 \( y = 4 \sin x + 3 \cos x \) 的最大值是\fillinblank{}.
\end{problem}

\begin{answer}
\(5\)
\end{answer}

\begin{solution}
由 \(\sin^2x+\cos^2x=1\)，有
\[
4\sin x+3\cos x=5\left(\frac45\sin x+\frac35\cos x\right)
\]
取满足 \(\cos\varphi=\frac45\)、\(\sin\varphi=\frac35\) 的角 \(\varphi\)，上式为 \(5\sin(x+\varphi)\). 因为 \(-1\le\sin(x+\varphi)\le1\)，其最大值为 \(5\)，且可以取得.
\end{solution}

\begin{problem}
复数 \( 2+3\i \)（\( \i \) 是虚数单位）的模是\fillinblank{}.
\end{problem}

\begin{answer}
\(\sqrt{13}\)
\end{answer}

\begin{solution}
复数 \(z=2+3\i\) 的模满足 \(|z|=\sqrt{(\operatorname{Re}z)^2+(\operatorname{Im}z)^2}\). 因此
\[
|2+3\i|=\sqrt{2^2+3^2}=\sqrt{13}
\]
\end{solution}

\begin{problem}
在 \( \triangle ABC \) 中，角 \(A\)，\(B\)，\(C\) 所对边长分别为 \(a\)，\(b\)，\(c\)，若 \(a = 5\)，\(b = 8\)，\(B = 60^\circ\)，则 \( b =\fillinblank{}\).
\end{problem}

\begin{answer}
\(8\)
\end{answer}

\begin{solution}
按原题题面，已知条件中直接给出边 \(b=8\)，所以填空处所求的 \(b\) 就是 \(8\). 原参考答案中的 \(7\) 与该题面不一致；\(7\) 是把题面改为 \(c=8\) 后由余弦定理求得的数值，不能作为原题 \(b\) 的答案.
\end{solution}

\begin{problem}
在如图所示的正方体 \( ABCD - A_{1}B_{1}C_{1}D_{1} \) 中，异面直线 \(A_{1}B\) 与 \(B_{1}C\) 所成角的大小为\fillinblank{}.

\bitmapfigure[max width=0.34\linewidth]{img/2013/shanghai_spring/q9_fig1.png}
\end{problem}

\begin{answer}
\(\frac{\pi}{3}\)
\end{answer}

\begin{solution}
设正方体棱长为 \(s\). 取坐标 \(A(0,0,0)\)、\(B(s,0,0)\)、\(C(s,s,0)\)、\(A_1(0,0,s)\)、\(B_1(s,0,s)\). 则
\(\overrightarrow{A_1B}=(s,0,-s)\)，\(\overrightarrow{B_1C}=(0,s,-s)\).
两向量的夹角就是两异面直线所成的锐角，且
\[
\cos\theta=\frac{\overrightarrow{A_1B}\cdot\overrightarrow{B_1C}}{|\overrightarrow{A_1B}|\,|\overrightarrow{B_1C}|}
=\frac{s^2}{(s\sqrt2)(s\sqrt2)}=\frac12
\]
由于 \(0<\theta<\frac\pi2\)，故 \(\theta=\frac\pi3\).
\end{solution}

\begin{problem}
从 \(4\) 名男同学和 \(6\) 名女同学中随机选取 \(3\) 人参加某社团活动，选出的 \(3\) 人中男女同学都有的概率为\fillinblank{}（结果用数值表示）.
\end{problem}

\begin{answer}
\(\frac{4}{5}\)
\end{answer}

\begin{solution}
从 \(10\) 人中任选 \(3\) 人，共有 \(\mathrm C_{10}^{3}\) 种等可能选法. 选出的 \(3\) 人男女都有的反面是全为男生或全为女生，分别有 \(\mathrm C_4^3\) 种和 \(\mathrm C_6^3\) 种，故
\[
P=1-\frac{\mathrm C_4^3+\mathrm C_6^3}{\mathrm C_{10}^3}
=1-\frac{4+20}{120}=\frac45
\]
\end{solution}

\begin{problem}
若等差数列的前 \(6\) 项和为 \(23\)，前 \(9\) 项和为 \(57\)，则数列的前 \(n\) 项和 \(S_{n}=\)\fillinblank{}.
\end{problem}

\begin{answer}
\(\frac{5n^{2}-7n}{6}\)
\end{answer}

\begin{solution}
设等差数列首项为 \(a_1\)，公差为 \(d\). 由
\(S_6=\frac{6}{2}(2a_1+5d)=23\)，\(S_9=\frac{9}{2}(2a_1+8d)=57\)
得 \(2a_1+5d=\frac{23}{3}\)、\(2a_1+8d=\frac{38}{3}\). 两式相减得 \(d=\frac53\)，代回得 \(a_1=-\frac13\). 因而
\[
S_n=\frac n2\left(2a_1+(n-1)d\right)
=\frac n2\left(-\frac23+\frac{5(n-1)}3\right)
=\frac{5n^2-7n}{6}
\]
\end{solution}

\begin{problem}
\(36\) 的所有正约数之和可按如下方法得到：因为 \( 36=2^{2}\times3^{2} \)，所以 \(36\) 的所有正约数之和为 \(\displaystyle (1+3+3^{2})+(2+2\times3+2\times3^{2})+(2^{2}+2^{2}\times3+2^{2}\times3^{2})=(1+2+2^{2})(1+3+3^{2})=91\). 参照上述方法，可求得 \(2000\) 的所有正约数之和为\fillinblank{}.
\end{problem}

\begin{answer}
\(4836\)
\end{answer}

\begin{solution}
分解 \(2000=2^4\times5^3\). 每个正约数唯一写成 \(2^i5^j\)，其中 \(0\le i\le4\)、\(0\le j\le3\)，所以所有正约数之和为
\[
(1+2+2^2+2^3+2^4)(1+5+5^2+5^3)=31\times156=4836
\]
\end{solution}

\section{单选题}

\begin{problem}
展开式为 \(ad-bc\) 的行列式是（\quad）

\choices
    {\( \begin{vmatrix} a & b \\ d & c \end{vmatrix} \)}
    {\( \begin{vmatrix} a & c \\ b & d \end{vmatrix} \)}
    {\( \begin{vmatrix} a & d \\ b & c \end{vmatrix} \)}
    {\( \begin{vmatrix} b & a \\ d & c \end{vmatrix} \)}
\end{problem}

\begin{answer}
B
\end{answer}

\begin{solution}
二阶行列式满足 \(\begin{vmatrix}u&v\\w&t\end{vmatrix}=ut-vw\). 四个选项的展开式中，只有
\[
\begin{vmatrix}a&c\\b&d\end{vmatrix}=ad-bc
\]
因此应选 B.
\end{solution}

\begin{problem}
设 \( f^{-1}(x) \) 为函数 \( f(x)=\sqrt{x} \) 的反函数，下列结论正确的是（\quad）

\choices
    {\( f^{-1}(2)=2 \)}
    {\( f^{-1}(2)=4 \)}
    {\( f^{-1}(4)=2 \)}
    {\( f^{-1}(4)=4 \)}
\end{problem}

\begin{answer}
B
\end{answer}

\begin{solution}
函数 \(f(x)=\sqrt{x}\) 的定义域和值域均为 \([0,+\infty)\). 由 \(y=\sqrt{x}\) 得 \(x=y^2\)，所以 \(f^{-1}(x)=x^2\)（\(x\ge0\)）. 因而 \(f^{-1}(2)=4\)，而 \(f^{-1}(4)=16\)，只有选项 B 正确.
\end{solution}

\begin{problem}
直线 \( 2x - 3y + 1 = 0 \) 的一个方向向量是（\quad）

\choices
    {\( (2, -3) \)}
    {\( (2, 3) \)}
    {\( (-3, 2) \)}
    {\( (3, 2) \)}
\end{problem}

\begin{answer}
D
\end{answer}

\begin{solution}
直线 \(2x-3y+1=0\) 的法向量为 \((2,-3)\). 设方向向量为 \((u,v)\)，则必须满足 \(2u-3v=0\). 检查四个选项，只有 \((3,2)\) 满足 \(2\times3-3\times2=0\)，所以应选 D.
\end{solution}

\begin{problem}
函数 \( f(x)=x^{-\frac{1}{2}} \) 的大致图像是（\quad）

\choices
    {\choicebitmap[width=0.15\paperwidth]{img/2013/shanghai_spring/q16_optA.png}}
    {\choicebitmap[width=0.15\paperwidth]{img/2013/shanghai_spring/q16_optB.png}}
    {\choicebitmap[width=0.15\paperwidth]{img/2013/shanghai_spring/q16_optC.png}}
    {\choicebitmap[width=0.15\paperwidth]{img/2013/shanghai_spring/q16_optD.png}}
\end{problem}

\begin{answer}
A
\end{answer}

\begin{solution}
函数 \(f(x)=x^{-1/2}=\frac1{\sqrt{x}}\) 的定义域为 \(x>0\)，函数值始终为正. 且
\(f'(x)=-\frac{1}{2x^{3/2}}<0\)，\((x>0)\)
所以图像在第一象限内从左上方向右下方递减. 四个图像中符合这一特征的是 A.
\end{solution}

\begin{problem}
如果 \(a < b < 0\)，那么下列不等式成立的是（\quad）

\choices
    {\( \frac{1}{a} < \frac{1}{b} \)}
    {\( ab < b^{2} \)}
    {\( -ab < -a^{2} \)}
    {\( -\frac{1}{a} < -\frac{1}{b} \)}
\end{problem}

\begin{answer}
D
\end{answer}

\begin{solution}
因为 \(a<b<0\)，函数 \(y=\frac1x\) 在负半轴上递减，所以 \(\frac1a>\frac1b\)，选项 A 不成立. 又因 \(b<0\)，由 \(a<b\) 得 \(ab>b^2\)，选项 B 不成立；由 \(a<0\) 且 \(b-a>0\) 得 \(a(b-a)<0\)，即 \(ab<a^2\)，所以 \(-ab>-a^2\)，选项 C 不成立. 将 \(\frac1a>\frac1b\) 两边乘以 \(-1\)，得 \(-\frac1a<-\frac1b\)，故应选 D.
\end{solution}

\begin{problem}
若复数 \( z_{1} \)，\( z_{2} \) 满足 \( z_{1}=\overline{z_{2}} \)，则 \( z_{1} \)，\( z_{2} \) 在复数平面上对应的点 \(Z_{1}\)，\(Z_{2}\)（\quad）

\choices
    {关于 \(x\) 轴对称}
    {关于 \(y\) 轴对称}
    {关于原点对称}
    {关于直线 \(y = x\) 对称}
\end{problem}

\begin{answer}
A
\end{answer}

\begin{solution}
设 \(z_2=u+v\i\)，则 \(z_1=\overline{z_2}=u-v\i\). 在复平面上，\(z_2\) 对应点的坐标为 \((u,v)\)，\(z_1\) 对应点的坐标为 \((u,-v)\)，两点关于 \(x\) 轴对称. 因此应选 A.
\end{solution}

\begin{problem}
\( (1+x)^{10} \) 的二项展开式中的一项是（\quad）

\choices
    {\(45x\)}
    {\( 90x^{2} \)}
    {\( 120x^{3} \)}
    {\( 252x^{4} \)}
\end{problem}

\begin{answer}
C
\end{answer}

\begin{solution}
由二项式定理，\(x^k\) 的系数为 \(\mathrm C_{10}^{k}\). 四个选项分别对应
\(45x\ne10x\)，\(90x^2\ne\mathrm C_{10}^{2}x^2=45x^2\)
\(\mathrm C_{10}^{3}x^3=120x^3\)，\(\mathrm C_{10}^{4}x^4=210x^4\ne252x^4\).
所以应选 C.
\end{solution}

\begin{problem}
既是偶函数又在区间 \( (0, \pi) \) 上单调递减的函数是（\quad）

\choices
    {\( y = \sin x \)}
    {\( y = \cos x \)}
    {\( y = \sin 2x \)}
    {\( y = \cos 2x \)}
\end{problem}

\begin{answer}
B
\end{answer}

\begin{solution}
\(y=\sin x\) 和 \(y=\sin2x\) 都是奇函数，不符合偶函数条件. \(y=\cos x\) 是偶函数，且在 \(0<x<\pi\) 时有 \(y'=-\sin x<0\)，因此单调递减. 虽然 \(y=\cos2x\) 也是偶函数，但其导数 \(-2\sin2x\) 在 \((0,\frac\pi2)\) 与 \((\frac\pi2,\pi)\) 上符号不同，不是整个区间上的单调递减函数. 故选 B.
\end{solution}

\begin{problem}
若两个球的表面积之比为 \(1:4\)，则这两个球的体积之比为（\quad）

\choices
    {\(1:2\)}
    {\(1:4\)}
    {\(1:8\)}
    {\(1:16\)}
\end{problem}

\begin{answer}
C
\end{answer}

\begin{solution}
设两个球的半径分别为 \(r_1\)，\(r_2\). 由表面积比得
\[
4\pi r_1^2:4\pi r_2^2=1:4
\]
所以 \(r_1:r_2=1:2\). 体积与半径的三次方成正比，故
\[
V_1:V_2=r_1^3:r_2^3=1:8
\]
因此应选 C.
\end{solution}

\begin{problem}
设全集 \( U = \R \)，下列集合运算结果为 \( \R \) 的是（\quad）

\choices
    {\( \Z \cup \complement_U \N \)}
    {\( \N \cap \complement_U \N \)}
    {\( \complement_U (\complement_U \emptyset) \)}
    {\( \complement_U \{0\} \)}
\end{problem}

\begin{answer}
A
\end{answer}

\begin{solution}
全集为 \(\R\). 对选项 A，任意实数若属于 \(\N\)，则也属于 \(\Z\)；若不属于 \(\N\)，则属于 \(\complement_{\R}\N\)，故
\[
\Z\cup\complement_{\R}\N=\R
\]
选项 B 的交集为空，选项 C 为 \(\complement_{\R}(\complement_{\R}\emptyset)=\emptyset\)，选项 D 为 \(\R\smallsetminus\{0\}\)，均不等于 \(\R\). 所以应选 A.
\end{solution}

\begin{problem}
已知 \( a \)，\( b \)，\( c \in \R \)，“\( b^2 - 4ac < 0 \)”是“函数 \( f(x) = ax^2 + bx + c \) 的图像恒在 \( x \) 轴上方”的（\quad）

\choices
    {充分非必要条件}
    {必要非充分条件}
    {充要条件}
    {既非充分又非必要条件}
\end{problem}

\begin{answer}
D
\end{answer}

\begin{solution}
按题面只给出 \(a\)，\(b\)，\(c\in\R\)，因此允许 \(a=0\). 先看充分性：取 \(a=-1\)，\(b=0\)，\(c=-1\)，则
\(b^2-4ac=-4<0\)，但 \(f(x)=-x^2-1<0\)，图像在 \(x\) 轴下方，故该条件不是充分条件.

再看必要性：取 \(a=0\)，\(b=0\)，\(c=1\)，则 \(f(x)=1\) 的图像恒在 \(x\) 轴上方，但 \(b^2-4ac=0\)，不满足小于零，故该条件也不是必要条件. 因而它既非充分又非必要，应选 D.
\end{solution}

\begin{problem}
已知 \(A\)，\(B\) 为平面内两定点，过该平面内动点 \(M\) 作直线 \(AB\) 的垂线，垂足为 \(N\). 若 \( \overrightarrow{MN}^{2} = \lambda \overrightarrow{AN} \cdot \overrightarrow{NB} \)，其中 \( \lambda \) 为常数，则动点 \(M\) 的轨迹不可能是（\quad）

\choices
    {圆}
    {椭圆}
    {抛物线}
    {双曲线}
\end{problem}

\begin{answer}
C
\end{answer}

\begin{solution}
设 \(AB=L>0\)，取 \(AB\) 为 \(x\) 轴，令 \(A=(0,0)\)、\(B=(L,0)\)、\(N=(t,0)\)、\(M=(t,y)\). 则
\(\overrightarrow{MN}=(0,-y)\)，\(\overrightarrow{AN}=(t,0)\)，\(\overrightarrow{NB}=(L-t,0)\)
从而轨迹满足
\[
y^2=\lambda t(L-t)
\]

当 \(\lambda>0\) 且 \(0\le t\le L\) 时，配方得
\[
\frac{(t-L/2)^2}{L^2/4}+\frac{y^2}{\lambda L^2/4}=1
\]
这是椭圆，且 \(\lambda=1\) 时为圆. 当 \(\lambda<0\) 时，令 \(\mu=-\lambda>0\)，则
\[
y^2=\mu t(t-L)
\]
即
\[
\frac{(t-L/2)^2}{L^2/4}-\frac{y^2}{\mu L^2/4}=1
\]
这是双曲线. 当 \(\lambda=0\) 时，轨迹退化为直线 \(y=0\). 因而无论哪种情况都不可能得到抛物线，所以应选 C.
\end{solution}

\section{解答题}

\begin{problem}
如图，在正三棱柱 \( ABC-A_{1}B_{1}C_{1} \) 中，\( AA_{1}=6 \)，异面直线 \(BC_{1}\) 与 \(AA_{1}\) 所成角的大小为 \( \frac{\pi}{6} \)，求该三棱柱的体积.

\bitmapfigure[max width=0.26\linewidth]{img/2013/shanghai_spring/q25_fig1.png}
\end{problem}

\begin{solution}
因为 \(CC_{1}\parallel AA_{1}\)，所以 \(\angle BC_{1}C\) 为异面直线 \(BC_{1}\) 与 \(AA_{1}\) 所成的角，即 \(\angle BC_{1}C=\frac{\pi}{6}\).
在直角三角形 \(BC_{1}C\) 中，
\[
BC=CC_{1}\tan\frac{\pi}{6}=6\cdot\frac{1}{\sqrt{3}}=2\sqrt{3}
\]
从而 \(S_{\triangle ABC}=\frac{\sqrt{3}}{4}\cdot BC^{2}=3\sqrt{3}\). 因此该三棱柱的体积为 \(V=S_{\triangle ABC}\cdot AA_{1}=3\sqrt{3}\cdot 6=18\sqrt{3}\).
\end{solution}

\begin{problem}
如图，某校有一块形如直角三角形 \(ABC\) 的空地，其中 \( \angle B \) 为直角，\(AB\) 长 \(40\text{ 米}\)，\(BC\) 长 \(50\text{ 米}\)，现欲在此空地上建造一间健身房，其占地形状为矩形，且 \(B\) 为矩形的一个顶点，求该健身房的最大占地面积.

\bitmapfigure[max width=0.22\linewidth]{img/2013/shanghai_spring/q26_fig1.png}
\end{problem}

\begin{solution}
设矩形为 \(EBFP\)，\(FP\) 长为 \(x\text{ 米}\)，其中 \(0<x<40\).
健身房占地面积为 \(y\text{ 平方米}\). 因为 \(\triangle CFP\sim\triangle CBA\)，\(\frac{FP}{BA}=\frac{CF}{CB}\)，即 \(\frac{x}{40}=\frac{50-BF}{50}\)，求得 \(BF=50-\frac{5}{4}x\).
从而
\[
y=BF\cdot FP=\left(50-\frac{5}{4}x\right)x=-\frac{5}{4}x^2+50x=-\frac{5}{4}(x-20)^2+500\le 500
\]
当且仅当 \(x=20\) 时，等号成立.
答：该健身房的最大占地面积为 \(500\text{ 平方米}\).
\end{solution}

\begin{problem}
已知数列 \( \{a_n\} \) 的前 \( n \) 项和为 \( S_n = -n^2 + n \)，数列 \( \{b_n\} \) 满足 \( b_n = 2^{a_n} \)，求 \( \lim\limits_{n \to \infty} (b_1 + b_2 + \cdots + b_n) \).
\end{problem}

\begin{solution}
当 \(n\ge 2\) 时，
\[
a_n=S_n-S_{n-1}=(-n^{2}+n)-\bigl(-(n-1)^{2}+(n-1)\bigr)=-2n+2
\]
且 \(a_1=S_1=0\). 因此 \(b_n=2^{a_n}=2^{-2n+2}=\bigl(\frac14\bigr)^{n-1}\). 所以数列 \(\{b_n\}\) 是首项为 \(1\)，公比为 \(\frac14\) 的无穷等比数列，
\(\lim\limits_{n\to\infty}(b_1+b_2+\cdots+b_n)=\frac{1}{1-\frac14}=\frac43\).
\end{solution}

\begin{problem}
已知椭圆 \(C\) 的两个焦点分别为 \( F_{1}(-1,0) \)，\( F_{2}(1,0) \)，短轴的两个端点分别为 \(B_{1}\)，\(B_{2}\).

\begin{enumerate}
\item 若 \( \triangle F_{1}B_{1}B_{2} \) 为等边三角形，求椭圆 \(C\) 的方程；
\item 若椭圆 \(C\) 的短轴长为 \(2\)，过点 \(F_{2}\) 的直线 \( l \) 与椭圆 \(C\) 相交于 \(P\)，\(Q\) 两点，且 \( \overrightarrow{F_{1}P} \perp \overrightarrow{F_{1}Q} \)，求直线 \( l \) 的方程.
\end{enumerate}
\end{problem}

\begin{solution}
\begin{enumerate}
\item 设椭圆 \(C\) 的方程为 \(\frac{x^2}{a^2}+\frac{y^2}{b^2}=1\)（\(a>b>0\)），则 \(c=1\)，\(a^2=b^2+c^2\). 由 \(\triangle F_1B_1B_2\) 为等边三角形，得 \(F_1B_1^2=1+b^2=(2b)^2\)，所以 \(b^2=\frac13\)，\(a^2=\frac43\). 故椭圆 \(C\) 的方程为
\(\frac{x^2}{4/3}+\frac{y^2}{1/3}=1\).

\item 短轴长为 \(2\)，故 \(b=1\)，\(a^2=2\)，椭圆方程为 \(\frac{x^2}{2}+y^2=1\). 当直线 \(l\) 的斜率不存在时，其方程为 \(x=1\)，交点为 \((1,1)\)、\((1,-1)\)，且
\(\overrightarrow{F_1P}=(2,1)\)，\(\overrightarrow{F_1Q}=(2,-1)\)，\(\overrightarrow{F_1P}\cdot\overrightarrow{F_1Q}=3\ne0\)
不符合垂直条件；当直线 \(l\) 的斜率存在时，设直线 \(l\) 的方程为 \(y=k(x-1)\). 联立椭圆方程得
\((2k^2+1)x^2-4k^2x+2(k^2-1)=0\).
设 \(P(x_1,y_1)\)，\(Q(x_2,y_2)\)，则
\(x_1+x_2=\frac{4k^2}{2k^2+1}\)，\(x_1x_2=\frac{2(k^2-1)}{2k^2+1}\)，
且 \(y_i=k(x_i-1)\). 由 \(\overrightarrow{F_1P}\perp\overrightarrow{F_1Q}\)，
\[
0=(x_1+1)(x_2+1)+y_1y_2
=(k^2+1)x_1x_2-(k^2-1)(x_1+x_2)+k^2+1
=\frac{7k^2-1}{2k^2+1}
\]
上述关于 \(x_1\)，\(x_2\) 的二次方程的判别式为
\[
\Delta=8(k^2+1)>0
\]
所以所得斜率对应两个不同的交点；故 \(k=\pm\frac{\sqrt7}{7}\) 均满足题设条件，直线 \(l\) 的方程为 \(x+\sqrt7y-1=0\) 或 \(x-\sqrt7y-1=0\).
\end{enumerate}
\end{solution}

\begin{problem}
已知抛物线 \( C: y^{2} = 4x \) 的焦点为 \(F\).

\begin{enumerate}
\item 点 \(A\)，\(P\) 满足 \( \overrightarrow{AP} = -2\overrightarrow{FA} \). 当点 \(A\) 在抛物线 \(C\) 上运动时，求动点 \(P\) 的轨迹方程；
\item 在 \(x\) 轴上是否存在点 \(Q\)，使得点 \(Q\) 关于直线 \( y = 2x \) 的对称点在抛物线 \(C\) 上？如果存在，求所有满足条件的点 \(Q\) 的坐标；如果不存在，请说明理由.
\end{enumerate}
\end{problem}

\begin{solution}
\begin{enumerate}
\item 设动点 \(P\) 的坐标为 \((x,y)\)，点 \(A\) 的坐标为 \((x_A,y_A)\). 由 \(F=(1,0)\) 及 \(\overrightarrow{AP}=-2\overrightarrow{FA}\)，得 \((x-x_A,y-y_A)=-2(x_A-1,y_A)\)，
从而 \(x_A=2-x\)，\(y_A=-y\). 又 \(A\) 在抛物线 \(C:y^2=4x\) 上，故 \(y_A^2=4x_A\)，即动点 \(P\) 的轨迹方程为 \(y^2=8-4x\).

\item 设点 \(Q\) 的坐标为 \((t,0)\)，对称点为 \(Q'(u,v)\). 因为 \(QQ'\perp y=2x\)，且 \(QQ'\) 的中点在 \(y=2x\) 上，所以
\((u-t)+2v=0\)，\(\frac v2=2\cdot\frac{t+u}{2}\).
解得 \(u=-\frac35t\)，\(v=\frac45t\)，即 \(Q'=\left(-\frac35t,\frac45t\right)\). 若 \(Q'\) 在 \(C\) 上，则 \(\left(\frac45t\right)^2=4\left(-\frac35t\right)\)，
解得 \(t=0\) 或 \(t=-\frac{15}{4}\). 所以存在满足题意的点 \(Q\)，其坐标为 \((0,0)\) 和 \(\left(-\frac{15}{4},0\right)\).
\end{enumerate}
\end{solution}

\begin{problem}
在平面直角坐标系 \(xOy\) 中，点 \(A\) 在 \(y\) 轴正半轴上，点 \( P_{n} \) 在 \(x\) 轴上，其横坐标为 \( x_{n} \)，且 \( \{x_{n}\} \) 是首项为 \(1\)，公比为 \(2\) 的等比数列，记 \( \angle P_{n}AP_{n+1} = \theta_{n} \)，\( n \in \N^{*} \).

\begin{enumerate}
\item 若 \( \theta_{3}=\arctan\frac{1}{3} \)，求点 \(A\) 的坐标；
\item 若点 \(A\) 的坐标为 \( (0, 8\sqrt{2}) \)，求 \( \theta_{n} \) 的最大值及相应 \(n\) 的值.
\end{enumerate}

\bitmapfigure[max width=0.55\linewidth]{img/2013/shanghai_spring/q30_fig1.png}
\end{problem}

\begin{solution}
\begin{enumerate}
\item 设 \(A(0,t)\). 根据题意，\(P_n=(2^{n-1},0)\). 由 \(\theta_3=\arctan\frac13\)，
\[
\tan\theta_3=\tan(\angle OAP_4-\angle OAP_3)=\frac{\frac{x_4}{t}-\frac{x_3}{t}}{1+\frac{x_4x_3}{t^2}}=\frac{4t}{t^2+32}=\frac13
\]
解得 \(t=4\) 或 \(t=8\)，故点 \(A\) 的坐标为 \((0,4)\) 或 \((0,8)\).

\item 由题意 \(A=(0,8\sqrt2)\)，且 \(\tan\angle OAP_n=\frac{2^{n-1}}{8\sqrt2}\). 因此
\[
\tan\theta_n=\frac{\frac{2^n}{8\sqrt2}-\frac{2^{n-1}}{8\sqrt2}}{1+\frac{2^n}{8\sqrt2}\cdot\frac{2^{n-1}}{8\sqrt2}}=\frac{1}{\frac{16\sqrt2}{2^n}+\frac{2^n}{8\sqrt2}}
\]
由均值不等式，
\[
\frac{16\sqrt2}{2^n}+\frac{2^n}{8\sqrt2}\ge 2\sqrt2
\]
且当且仅当 \(n=4\) 时取等号. 由于 \(0<\theta_n<\frac\pi2\)，故 \(\theta_n\) 的最大值为 \(\arctan\frac{\sqrt2}{4}\)，相应 \(n=4\).
\end{enumerate}
\end{solution}

\begin{problem}
已知真命题：“函数 \( y = f(x) \) 的图像关于点 \( P(a, b) \) 成中心对称图形”的充要条件为“函数 \( y = f(x + a) - b \) 是奇函数”.

\begin{enumerate}
\item 将函数 \( g(x)=x^{3}-3x^{2} \) 的图像向左平移 \(1\) 个单位，再向上平移 \(2\) 个单位，求此时图像对应的函数解析式，并利用题设中的真命题求函数 \( g(x) \) 图像对称中心的坐标；
\item 求函数 \( h(x)=\log_{2}\frac{2x}{4-x} \) 图像对称中心的坐标；
\item 已知命题：“函数 \( y = f(x) \) 的图像关于某直线成轴对称图像”的充要条件为“存在实数 \(a\) 和 \(b\)，使得函数 \( y = f(x + a) - b \) 是偶函数”. 判断该命题的真假. 如果是真命题，请给予证明；如果是假命题，请说明理由，并类比题设的真命题对它进行修改，使之成为真命题（不必证明）.
\end{enumerate}
\end{problem}

\begin{solution}
\begin{enumerate}
\item 平移后图像对应的函数解析式为 \(y=(x+1)^3-3(x+1)^2+2=x^3-3x\).
由于函数 \(y=x^3-3x\) 是奇函数，由题设真命题知，函数 \(g(x)\) 图像对称中心的坐标是 \((1,-2)\).

\item 设 \(h(x)=\log_2\frac{2x}{4-x}\) 的对称中心为 \(P(a,b)\). 由题设知函数 \(h(x+a)-b\) 是奇函数. 令 \(f(x)=h(x+a)-b\)，则 \(f(x)=\log_2\frac{2x+2a}{4-a-x}-b\).
由不等式 \(\frac{2x+2a}{4-a-x}>0\) 的解集为 \((-a,4-a)\)，其区间中点为 \(2-a\). 因 \(f\) 为奇函数，其定义域须关于原点对称，故 \(2-a=0\)，得 \(a=2\). 此时 \(f(x)=\log_2\frac{2(x+2)}{2-x}-b\)，\(x\in(-2,2)\).
任取 \(x\in(-2,2)\)，有
\[
f(-x)+f(x)
=\log_2\left(\frac{2(2-x)}{2+x}\cdot\frac{2(x+2)}{2-x}\right)-2b
=\log_2 4-2b
=2-2b
\]
由 \(f\) 为奇函数，\(f(-x)+f(x)=0\)，故 \(b=1\). 所以函数 \(h(x)\) 图像对称中心的坐标是 \((2,1)\).

\item 该命题是假命题. 举反例说明：函数 \(f(x)=x\) 的图像关于直线 \(y=x\) 成轴对称图像，但是对任意实数 \(a\) 和 \(b\)，函数 \(y=f(x+a)-b\)，即 \(y=x+a-b\) 总不是偶函数.

修改后的真命题：“函数 \(y=f(x)\) 的图像关于直线 \(x=a\) 成轴对称图像”的充要条件是“函数 \(y=f(x+a)\) 是偶函数”.
\end{enumerate}
\end{solution}
